\documentclass[12pt]{article}

\usepackage{graphicx}
\usepackage{amssymb}
\usepackage{amsmath}
\usepackage{mathtools}
\usepackage{authblk}
\usepackage{titling}
\usepackage{hyperref}
\usepackage[final]{microtype}

\title{Operational Schwarzschild Horizons\\ from Hamiltonian Quantum Mechanics}
\author{Sergei Slobodov\thanks{sergei@slobodov.com}}
\date{\today}

\textwidth = 6.5 in
\textheight = 9.0 in
\oddsidemargin = 0.0 in
\evensidemargin = 0.0 in
\topmargin = 0.0 in
\headheight = 0.0 in
\headsep = 0.0 in
\parskip = 0.2in
\parindent = 0.0in
\emergencystretch=1em

\hypersetup{colorlinks=true, linkcolor=blue, citecolor=blue, urlcolor=blue}

\newtheorem{theorem}{Theorem}
\newtheorem{corollary}[theorem]{Corollary}
\newtheorem{definition}{Definition}

\newcommand{\be}{\begin{equation}}
\newcommand{\ee}{\end{equation}}
\newcommand{\bea}{\begin{eqnarray}}
\newcommand{\eea}{\end{eqnarray}}
\newcommand{\Tr}{\mathrm{Tr}}
\newcommand{\Smicro}{S_{\rm micro}}
\newcommand{\SBH}{S_{\rm BH}}
\newcommand{\HB}{\mathcal H_{\rm B}}
\newcommand{\HR}{\mathcal H_{\rm R}}
\newcommand{\ket}[1]{|#1\rangle}
\newcommand{\bra}[1]{\langle #1|}

\renewcommand\Authands{ and }

\begin{document}

\begin{titlingpage}
\pagenumbering{gobble}
\maketitle

\begin{abstract}
We introduce a class of Hamiltonian models for the exterior, operational
properties of Schwarzschild evaporation.  The models are quantum systems with
an energy-resolved black-hole Hilbert space, weak emission into outgoing
radiation modes, and mixing within fixed energy shells.  We impose
three inputs: the Schwarzschild density of states, boundary accessibility of
the leading entropy through area-many emission operators, and decoupling of the
induced shrinking emission channel.  Given these inputs, the models reproduce
the Hawking spectrum, negative heat capacity, Schwarzschild flux and lifetime
scalings, the Page curve, the island entropy formula, complementary recovery,
and state-dependent mirror operators.  Conversely, within this model class, the
three inputs are calibrated by the corresponding exterior phenomenology: the
density of states by the temperature law, the boundary emission algebra by the
luminosity law, and decoupling by complementary recovery.  Thus the operational
horizon provides a Hamiltonian separation between the exterior information
structure of Schwarzschild evaporation and the microscopic origin of the
black-hole state count, boundary emission algebra, and mixing dynamics.
\end{abstract}
\end{titlingpage}
\pagenumbering{arabic}

\section{Introduction}

Semiclassical black holes possess an exterior operational layer visible in
their radiation and recovery properties: Hawking thermality, area-law emission
and absorption, the Page transition, island reconstruction, complementarity,
and a decoding barrier
\cite{Hawking1975,Page1993Average,Penington2019,HaydenPreskill2007,
HarlowHayden2013}.  The purpose of this paper is to isolate that operational
layer in a model where every quantum mechanical step is explicit.

\begin{definition}[Operational Schwarzschild horizon]
An operational Schwarzschild horizon is the exterior signature package
consisting of detailed-balance thermality at $T\propto M^{-1}$, emission
strength proportional to area, the Page/island transition of the radiation
entropy, complementary recovery of code information across that transition,
state-dependent mirror operators, and the associated complexity barrier.
\end{definition}

The central claim is that this package follows from ordinary Hamiltonian
quantum mechanics once three inputs are provided.

The inputs are:
\begin{itemize}
\item
the Schwarzschild microcanonical entropy
\be
  \Smicro(E)=cE^2 ,
\ee
\item
boundary accessibility of the leading entropy through an area-sized set of
weak emission operators, and
\item
decoupling, or equivalently typical encoding, for the shrinking emission
channel generated by the shell dynamics.
\end{itemize}

\begin{theorem}[Hamiltonian realization]
In the Hamiltonian model defined below, the three inputs above realize an
operational Schwarzschild horizon.  The density of states gives
detailed-balance Hawking thermality and negative heat capacity, boundary
accessibility gives the Schwarzschild emission strength, and decoupling gives
the Page/island entropy formula, complementary recovery, and the corresponding
state-dependent mirrors.  The computational barrier is inherited under the
standard Harlow-Hayden/Python's-lunch assumptions.
\end{theorem}

The first two inputs are thermodynamic.  Their necessity is elementary but
important: the density of states fixes the Hawking temperature, while the
area-sized emission algebra fixes the Schwarzschild luminosity.  The third
input is information-theoretic.  It is equivalent, in the standard
information--disturbance sense, to complementary recovery of code information
from the radiation after the Page transition.

This framing is close in spirit to the tensor-network models of Leutheusser
and Van Raamsdonk \cite{LeutheusserVanRaamsdonk2017}.  Those models isolate the
information-theoretic structure of unitary black-hole evaporation in terms of a
sequence of isometries.  The present model adds the energy, density of states,
weak-emission Hamiltonian, and boundary emission algebra needed to recover the
Schwarzschild thermodynamic scalings.  The conditional statement identifies
the state count, boundary accessibility, and decoupling as the microscopic
ingredients whose origin a completion must supply.

Section 2 defines the Hamiltonian model.  Section 3 derives the thermodynamic
relations and the necessity of the first two inputs.  Section 4 defines the
operational horizon package and shows how it follows from decoupling.  Section
5 states the converse recovery theorem and the resulting necessity trinity.
Section 6 formulates the geometric-completion and microscopic-origin
questions.  Technical details are collected in appendices.

\subsubsection*{Relation to earlier work}

Toy models of black-hole evaporation have long been used to isolate the
quantum-information part of the information problem
\cite{Avery2013,OsugaPage2018,PiroliSunderhaufQi2020}.  Tensor-network and
random-isometry models make Page-like information flow transparent, but usually
take the radiation step as kinematic data.  The model below keeps the same
isometric evaporation structure while deriving the one-step map from a
weak-emission Hamiltonian with an energy-dependent density of states.

The recovery layer also overlaps standard quantum-information and
cryptographic language.  Small radiation fragments are unauthorized for the
diary in the same sense as unauthorized sets in quantum secret sharing, and
the post-Page Hayden--Preskill transition is a decoupling/recovery threshold
rather than a new access-structure mechanism
\cite{CleveGottesmanLo1999,TerhalDiVincenzoLeung2001,
DiVincenzoHaydenTerhal2004}.  The point here is to identify which
thermodynamic and emission-algebra inputs make that standard recovery
structure look like Schwarzschild evaporation.

Island calculations and replica wormholes provide the gravitational derivation
of the Page curve \cite{Penington2019,AlmheiriEngelhardtMarolfMaxfield2019,
PeningtonShenkerStanfordYang2022}.  Here the corresponding entropy formula is
obtained as the Page/isometry result for a shrinking Hilbert space, with the
remaining microcanonical entropy replacing the area term.  This gives a
Hamiltonian version of the island entropy rule once the black-hole state count
and decoupling input are supplied.

\section{Hamiltonian models of the operational horizon}

The black-hole subsystem is a direct sum of energy shells
\be
  \HB=\bigoplus_E \mathcal H_E ,
  \qquad
  \dim\mathcal H_E=\exp[\Smicro(E)] .
\ee
For a four-dimensional Schwarzschild black hole we take
\be
  \Smicro(E)=cE^2 ,
  \qquad
  \beta(E)=\frac{d\Smicro}{dE}=2cE .
\ee
The outgoing radiation is a Fock space with modes labelled by frequency
$\omega$ and channel $\lambda$.

The Hamiltonian has the form
\be
  H=H_{\rm B}+H_{\rm R}+H_I+H_{\rm mix}.
\ee
Here $H_{\rm B}$ resolves the shell energy, $H_{\rm R}$ is the free radiation
Hamiltonian, $H_{\rm mix}$ mixes states within a fixed shell, and the
interaction
\be
  H_I=
  \sum_{\lambda,\mu}\int d\omega\,
  g_{E;\omega\lambda\mu}\,
  b^\dagger_{\omega\lambda} O_{E;\omega\lambda\mu}
  +{\rm h.c.}
\ee
connects $\mathcal H_E$ to $\mathcal H_{E-\omega}$.  The index $\mu$ labels
boundary-accessible emitters.  The physical radiation record carries
$(\omega,\lambda)$ while distinct emitters feed the same outgoing mode, so the
resolvable jump operator is the collective sum
\be
  K^{\rm coll}_{E;\omega\lambda}
  =
  \sum_\mu K_{E;\omega\lambda\mu}.
\ee
Random phases or ETH matrix elements make rates add incoherently while
preserving the single-branch isometric structure.

At Schwarzschild frequencies the emitted wavelength is of order the horizon
radius, so the asymptotic field carries only order-one angular channel capacity
per coherence time \cite{Page1976,BekensteinMayo2001,Pendry1983}.  The
area-sized count below is therefore a source-side count in the emission
algebra: it adds to the inclusive rate, while the physical radiation record is
the serial record of frequencies, outgoing channels, and time bins.  Full
exterior access to deposited information is controlled by this record together
with the shell mixing.

Boundary accessibility means that the number of comparably coupled active
emitters scales as the leading entropy,
\be
  N(E)\sim \Smicro(E)\sim E^2 .
\ee
This is the Hamiltonian version of the statement that the degrees carrying the
leading entropy are accessible to the emission algebra.

\section{Thermodynamic rigidity}

The density of states fixes the local emission spectrum.  For a transition
$E\to E-\omega$,
\be
  \frac{\rho(E-\omega)}{\rho(E)}
  =
  \exp[\Smicro(E-\omega)-\Smicro(E)]
  =
  e^{-\beta(E)\omega+c\omega^2}.
\ee
For $\omega/E\ll 1$ this is the Hawking Boltzmann factor at temperature
$T(E)=1/\beta(E)$, with the leading finite-energy correction
\cite{ParikhWilczek2000}.

Conversely, any microcanonical Hamiltonian model whose weak-emission spectrum
has the Schwarzschild temperature law must satisfy
\be
  \frac{dS}{dE}\propto E,
  \qquad
  S(E)\propto E^2
\ee
up to an additive constant.  Thus the state-count input is equivalent, within
this model class, to the Hawking temperature law.

Within this emission class, the luminosity calibrates the second input.  In
$3+1$ dimensions the density of radiation phase space gives a power
\be
  P(E)\sim N(E) T(E)^4 .
\ee
Schwarzschild evaporation has $P(E)\sim E^{-2}$.  Since $T(E)\sim E^{-1}$,
this requires
\be
  N(E)\sim E^2\sim \Smicro(E).
\ee
Thus the area-sized boundary emission algebra is the source-side realization
of the luminosity law.

Under one further assumption these two facts imply a corpuscular kinematics.
If the boundary-accessible cells are comparable energy-carrying constituents
rather than degeneracy labels or edge modes, then
\be
  \epsilon_{\rm cell}\sim \frac{E}{N(E)}\sim \frac{1}{E}\sim T(E).
\ee
The thermodynamic inputs then force a large-$N$ system of $N\sim E^2$ soft
quanta, each carrying order-one entropy and energy of order the temperature.
This is the kinematic content of the quantum $N$-portrait of Dvali and Gomez
\cite{DvaliGomez2013}, obtained here as a consequence of the two measured
scalings plus the energy-sharing assumption.

\section{The operational horizon package}

Let a narrow evaporation trajectory be represented by an isometry
\be
  V_E:\mathcal H_E\to
  \bigoplus_{\omega,\lambda}
  \mathcal H_{E-\omega}\otimes
  \ket{\omega,\lambda}_{\rm R}.
\ee
Composing these maps gives a shrinking-channel description of the evaporation
history.  If the emitted record has effective dimension
$D_{\rm R}(E)=\exp[\Delta S_{\rm rad}(E)]$ and the remaining shell has
$D_{\rm B}(E)=\exp[\Smicro(E)]$, decoupling gives
\be
  S(R;E)=
  \min\{\Delta S_{\rm rad}(E),\Smicro(E)\}+O(1).
\ee
This is the Page curve \cite{Page1993Average}.  At second Renyi order the
same calculation has two
contractions,
\be
  \Tr \rho_R^2
  \simeq
  e^{-\Delta S_{\rm rad}(E)}+
  e^{-\Smicro(E)},
\ee
which is the island min formula with the remaining microcanonical entropy
playing the role of the area term.

The operational layer is the recovery statement.  Before the Page transition,
small code subspaces are recoverable from the remaining black-hole subsystem
and decoupled from the radiation.  After the transition the assignment is
reversed: the same code information is recoverable from the radiation record
\cite{HaydenPreskill2007}.
This is the algebraic content of complementarity in the model.

For each fresh Hawking mode, a mirror operator exists on the purifier of that
mode.  Before the Page time the purifier is mostly the remaining subsystem;
after the Page time it is mostly early radiation.  The mirror is necessarily
state-dependent by the usual finite-dimensional counting obstruction to a
single linear operator satisfying the mirror relations on typical microstates
\cite{PapadodimasRaju2013,PapadodimasRaju2014}.
When the mirror is radiation-supported, decoding it is exponentially complex
under the standard Harlow-Hayden/Python's-lunch assumptions
\cite{HarlowHayden2013,BrownGharibyanPeningtonSusskind2020}.

\section{Necessity of the three inputs}

The previous sections show that the three inputs are sufficient.  In the model
class they are also calibrated by the three exterior layers they produce.

The state-count input is calibrated by the temperature law:
\be
  T(E)^{-1}=\frac{dS}{dE}.
\ee
Schwarzschild temperature therefore forces $S(E)\propto E^2$.

The boundary-accessibility input is calibrated by the luminosity law:
\be
  P(E)\sim N(E)T(E)^4.
\ee
Schwarzschild power therefore forces $N(E)\sim E^2$.

The decoupling input is calibrated by the recovery layer.  By the standard
information--disturbance theorem, recovery of a code from one subsystem is
equivalent, up to the usual error degradation, to decoupling of the
complementary subsystem from the reference
\cite{Schumacher1996,KretschmannSchlingemannWerner2008}.  Thus complementary
recovery across the Page transition is the operational measurement of the
decoupling assumption.

\begin{theorem}[Necessity within the model class]
Consider weak-emission microcanonical Hamiltonian models with a smooth
shrinking channel, $3+1$ dimensional radiation phase space, and comparable
boundary emission channels.  If such a model has the Schwarzschild temperature
law, the Schwarzschild luminosity law, and complementary recovery of code
information across the Page transition, then it must have
$S(E)=cE^2+O(E^0)$, an area-sized accessible emission algebra
$N(E)\sim E^2$, and decoupling of the complementary channel with the standard
information--disturbance error conversion.
\end{theorem}

This is the necessity trinity.  The three inputs are fixed by the exterior
Schwarzschild phenomenology they reproduce: temperature, luminosity, and
complementary recovery.

\section{Geometric completion}

The model realizes the operational horizon.  A geometric completion would
identify the algebraic structures constructed here with local near-horizon
physics in a Schwarzschild spacetime.  In particular, the partner operators
constructed here reproduce the algebraic mirror relations for fresh modes; a
semiclassical completion would identify such relations with local
horizon-crossing correlators.

The microscopic origin of the inputs is the other completion problem.  At fixed
local dimension and range, ordinary local finite-density systems obey
extensive state-count bounds rather than $S(E)\propto E^2$; this is the
holographic entropy mismatch
\cite{Bekenstein1981,tHooft1993,Susskind1995,Bousso2002,
CohenKaplanNelson1999,Yurtsever2003}.  A microscopic completion must explain the
super-Hagedorn state count, why the leading entropy is boundary-accessible,
and why the shell dynamics decouples the shrinking emission channel.

There is also a factorization assumption.  The model treats the black-hole
subsystem and the radiation as tensor factors, so information resides in the
black-hole subsystem until it is emitted.  In gravity, holography-of-information
arguments based on the Gauss law challenge this bookkeeping by suggesting that
bulk information is already available in asymptotic observables
\cite{Raju2022}.  The present model provides the clean factorized benchmark
against which Gauss-law modifications can be tested.

\section{Discussion}

In this paper, we have introduced Hamiltonian models for the exterior
operational properties of Schwarzschild evaporation.  Given a Schwarzschild
state count, a boundary-accessible emission algebra, and decoupling dynamics,
the models reproduce the thermal spectrum, flux scaling, Page/island entropy
rule, complementary recovery, and state-dependent mirror operators.  Within
the same model class, these inputs are calibrated by the temperature law, the
luminosity law, and the recovery property.

It would be interesting to investigate microscopic models in which the inputs
arise dynamically.  One direction is to derive the thermal tie between
scrambling and temperature from the emission interaction itself.  Virtual
emission and reabsorption can generate radiation-mediated couplings between
boundary emitters.  A full-rank operator-growth kernel with scale set by the
emission rate would give the fast-scrambling estimate
\cite{SekinoSusskind2008} from the same Hamiltonian that radiates.  A low-rank
kernel would sharpen the status of mixing as an independent input.

A second direction is to compare fluctuations.  The mean island formula has a
Hamiltonian permutation expansion matching the replica-wormhole saddle sum
after the identification $\SBH=\log d_B$.  Extending this comparison to the
variance would test whether second-order freeness is the microscopic avatar of
half-wormhole corrections.  These directions use the present models as a
controlled setting for adding, one at a time, further features of black-hole
physics.

\appendix

\section{Collective jumps and the resolvable record}

The physical radiation record labels frequency, channel, and time bin.  The
boundary emitter index is summed into the shared outgoing mode, so the
resolvable jump is collective:
\be
  K_m=\sum_\mu K_{m\mu}.
\ee
If the microscopic ETH variables are zero mean and have second and fourth
moments diagonal in the joint label $(m,\mu)$, then the collective matrix
element is again ETH:
\be
  \bra{\bar\nu}K_m\ket{\nu}
  =
  D^{-1/2} F_m(\bar E,\omega) R^{(m)}_{\bar\nu\nu},
  \qquad
  |F_m|^2=\sum_\mu |f_{m\mu}|^2 .
\ee
Rates therefore add while amplitudes have zero coherent mean.  The Dicke
$N^2$ enhancement is absent because the mean amplitude vanishes.  Fourth
moments inherit Wick factorization, with connected corrections suppressed by
the same ETH factors and by the additional averaging over emitters.  The
normalization condition holds up to cross terms with $\mu\ne\mu'$, which vanish
on shell average and enter the same branch-weight error budget as the
golden-rule approximation.

\section{Decoupling estimate}

Let $Q$ purify a code subspace of the shell and let $B$ denote the remaining
black-hole factor after one or more emission steps.  A standard second-moment
decoupling estimate
\cite{HaydenPreskill2007,HaydenHorodeckiWinterYard2008,Dupuis2014} gives
\be
  \left\|\rho_{QB}-\pi_Q\otimes\rho_B\right\|_1^2
  \lesssim
  \exp\{k+\Smicro(E)-\Delta S_{\rm rad}(E)\}
\ee
up to the ETH and golden-rule errors.  Thus a diary of entropy $k$ is
recoverable from the radiation after
\be
  \Delta S_{\rm rad}(E)-\Smicro(E)\gtrsim k .
\ee
The converse direction is the standard information--disturbance theorem:
high-fidelity recovery from radiation implies decoupling of the complementary
black-hole factor from the reference.

\bibliographystyle{unsrturl}
\bibliography{refs}

\end{document}
